\section{Conclusion}\label{sec:conclusion}
The linear bounds arise from a way of comparing bounded-degree algebraic
calculations across universe sizes. For a fixed uniform request over a
field of positive characteristic, the construction transports the actual
axiom, multiplication and target maps. It preserves the linear-space
operations defining the consequence spaces, including the cancellations
which make ordinary PC stronger at a given degree. A separate detection
argument reflects whether the prescribed target belongs to that space.
Together these steps make the construction apply to uniform consequences,
with refutation as the target-$1$ special case.

The quantitative gain is a linear onset and a linearly bounded prime-power
period for each degree threshold. Repeated lifting keeps the numerical
degree fixed while the universe grows. A nonrecurring degree can therefore
occur only before its own onset; the period bound also locates recurring
large degrees at proportionate universe sizes. This gives bounded degree
in $\mathrm T_1$, linear hardness at finite values in $\mathrm T_2$, linear
hardness on the hard residues in $\mathrm T_3$, and finite linear peaks
with a finite liminf in $\mathrm T_4$. Underivable requests retain degree
$\infty$, and the stated initial exceptions remain part of the results.

The linear degree lower bound is sharp in general. The injection example has matching
$\Theta(n)$ degrees for both systems, and the weak-bijection example has
finite $\Theta(n)$ peaks; its degree-one target extension has the same
peak behaviour. These witnesses establish optimality of the general
degree lower bound, while the theorem leaves the individual upper bounds to the
particular uniform request.

The uniform NS--PC comparison shows that these patterns belong to the
request itself: both systems assign it the same type, even when their
minimum degrees differ. Their common hardness spectrum gives the same
four alternatives as topological shapes in $\mathbb Z_p$, for all
uniform requests. With full explicit coefficient lists, including
zeros, the degree results give the stated polynomial upper bounds and
exponential size lower bounds. The first-order refutation subcase connects
this algebraic classification to logic: explicit reductions separate
any two types under the no-finite-model promise already stated in
LICS~2008~\cite[Theorem~1]{Riis}, refining its every-union result.
Thus the stronger quantitative bounds preserve the four-way
structure of the earlier classification while extending its algebraic
scope.

\paragraph{Proof sources and verification.}
The archive \texttt{degree-profiles-lean-source.zip} supplies the
curated proof sources, pinned Lean~4.26.0 and mathlib dependency data,
and the source, statement, axiom and import-closure audit commands
and their recorded execution provenance. The core degree, comparison,
classification and undecidability endpoints have Lean proofs; the
archive identifies their exact types and assumptions. The injection
upper certificate, the optimality deductions and the size presentation
are written arguments. The topological corollary in
Section~\ref{sec:model-theory} is likewise a written deduction from
periodicity, comparison and compactness of $\mathbb Z_p$; it is not
an addition to the Lean formalisation. Recorded fresh replay covers the named import
closures, with separately compiled scope adapters distinguished in the
evidence; it is not a fresh replay of every package module.

\paragraph{Acknowledgements.}
I used the AI systems Astra and Claude for sustained proof search,
formalisation, verification support and development of the exposition,
under my mathematical direction and editorial supervision. The proof
search built on my 53-page 2009 manuscript~\cite{RiisLong} and my
unpublished notes on the problem. I supplied these sources, set the
research programme, and checked and revised the arguments and their
presentation. The human work in proof complexity,
model theory, modular representation theory and algebraic combinatorics
cited throughout the paper provides important antecedents and connections,
including~\cite{NSFoundations,PCFoundations,Ajtai,Krajicek,KrajEuler,
BussDesign,BeameRiis,RiisSitharamTautologies,RiisSitharamModules,
CEF,Harman,HKPSYZ,BDM}. Responsibility
for the statements, interpretation, attribution and final text remains mine.
